\section{The critical inverse-\texorpdfstring{\(J\)}{J} interface}
\label{sec:critical-interface}

Let \(\tau_j\to0\) be a target angle scale.  The TPC application
will set \(\tau_j=J_j^{-1}\).

\begin{theorem}[Safe and defective transfer dichotomy]
\label{thm:critical-dichotomy}
Let \(P_j\) be orthogonal projections on prescribed finite frames,
let
\[
 B_j=P_jA_j\ne0,\qquad
 R_j=(I-P_j)A_j,
\]
and put
\[
 \kappa_j
 =
 \frac{\|(I-P_j)\one\|_2}{\sqrt{S_j}},
\qquad
 r_j=\frac{\|R_j\|_2}{\|B_j\|_2}.
\]
Suppose
\[
 \rho(B_j)\le C\tau_j.
\]
\begin{enumerate}[label=(\roman*),leftmargin=2.2em]
\item If \(\kappa_j=0\), then
  \[
  \rho(A_j)\le C\tau_j
  \]
  for arbitrary \(r_j\).
\item If
  \[
  \kappa_jr_j\le K\tau_j,
  \]
  then
  \[
  \rho(A_j)\le(C+K)\tau_j.
  \]
\end{enumerate}
\end{theorem}

\begin{proof}
Part (i) is \cref{cor:no-accuracy-certificate}.  For part (ii),
\cref{thm:defective-envelope} gives
\[
 \rho(A_j)
 \le
 \frac{\rho(B_j)+\kappa_jr_j}{\sqrt{1+r_j^2}}
 \le(C+K)\tau_j.
\]
\end{proof}

\begin{theorem}[Sharpness at the target scale]
\label{thm:critical-sharpness}
Uniform transfer of a zero-mode error target
\[
 |Z(A)-Z(PA)|
 \le C\tau\sqrt S\,\|PA\|_2
\]
from only the two norm parameters \(\kappa(P)\) and
\(r=\|(I-P)A\|_2/\|PA\|_2\) requires
\[
 \kappa(P)r=O(\tau)
\]
in order.  In particular, on any class where
\(\kappa(P)\ge\kappa_0>0\), the structure-free relative residual
requirement returns as
\[
 r=O(\tau).
\]
For \(\tau=J^{-1}\), this is the TPC-89 critical order.
\end{theorem}

\begin{proof}
Sufficiency is \cref{cor:product-criterion}.  For necessity, when
\(q=(I-P)\one\ne0\), choose \(R\) parallel to \(q\).
\Cref{thm:sharp-zero-defect} then gives the exact normalized defect
\[
 \frac{|Z(A)-Z(PA)|}{\sqrt S\,\|PA\|_2}
 =\kappa(P)r.
\]
No smaller uniform order follows from the two norms.  If
\(\kappa(P)\ge\kappa_0\), the stated residual order is necessary.
\end{proof}

\begin{example}[Sharp angle failure beyond \(J^{-1}\)]
\label{ex:J-failure}
Use the projection from
\cref{ex:unsafe-false-positive}, so \(\kappa=1\), and let
\[
 B_J=f,\qquad
 A_J=f+r_Je.
\]
Then
\[
 \rho(B_J)=0,\qquad
 \rho(A_J)=\frac{r_J}{\sqrt{1+r_J^2}}.
\]
If \(Jr_J\) is unbounded, then \(J\rho(A_J)\) is unbounded along a
subsequence.  Thus even a perfect projected zero mode need not
transfer the critical angle when the missing-constant residual is
larger than order \(J^{-1}\).
\end{example}

\begin{remark}[The product tradeoff]
\label{rem:product-tradeoff}
The defective case does not always require
\(r=O(J^{-1})\).  An approximately constant-preserving projection
with \(\kappa=o(1)\) can tolerate a larger residual, provided
\[
 \kappa r=O(J^{-1}).
\]
The two factors have different meanings: \(\kappa\) is an operator
defect known before seeing \(A\), while \(r\) measures the discarded
energy of the particular coefficient.  They should be recorded
separately before their product is compared with the target.
\end{remark}

\subsection{Unnormalized critical scale}

The same condition appears in the physical zero-mode normalization.
Assume along a critical family that
\begin{equation}
 S_X\le X^{o(1)}Q_X,\qquad
 D(B_X)\le X^{o(1)}Q_X^3J_X^2,
\label{eq:critical-upper-scales}
\end{equation}
where \(B_X=P_XA_X\ne0\).

\begin{proposition}[Defect at the \(Q^2\) scale]
\label{prop:Q2-defect}
Under \eqref{eq:critical-upper-scales}, put
\[
 \kappa_X
 =
 \frac{\|(I-P_X)\one\|_2}{\sqrt{S_X}},
\qquad
 r_X
 =
 \frac{\|(I-P_X)A_X\|_2}{\|P_XA_X\|_2}.
\]
Then
\begin{equation}
 |Z(A_X)-Z(B_X)|
 \le
 X^{o(1)}\kappa_Xr_X Q_X^2J_X.
\label{eq:Q2-defect}
\end{equation}
Hence a sufficient condition for
\[
 |Z(A_X)-Z(B_X)|
 \le X^{-\eta_Z+o(1)}Q_X^2
\]
is
\begin{equation}
 \boxed{
 \kappa_Xr_X
 \le
 \frac{X^{-\eta_Z+o(1)}}{J_X}.}
\label{eq:physical-product-target}
\end{equation}
For a constant-preserving projection, the defect is identically
zero and this approximation condition disappears.
\end{proposition}

\begin{proof}
By \cref{thm:sharp-zero-defect},
\[
 |Z(A_X)-Z(B_X)|
 \le
 \sqrt{S_X}\,\kappa_Xr_X\sqrt{D(B_X)}.
\]
Insert \eqref{eq:critical-upper-scales} to obtain
\eqref{eq:Q2-defect}; the final statement follows immediately.
\end{proof}

\begin{remark}[Zero-mode and energy goals remain distinct]
\label{rem:zero-energy-distinct}
Exact zero-mode preservation does not by itself prove that the
projected energy is large.  Conversely, a large projected energy
does not estimate its zero mode.  \Cref{prop:safe-dominance-ledger}
shows that proved projected bounds transfer in the correct
directions, but both inputs still require their own arguments.
\end{remark}
